% ===================================================================
% CHAPTER 7: EXPERIMENTS AND RESULTS
% ===================================================================
\chapter{Експерименти и резултати}
\label{chap:results}

\section{Експериментална поставеност}

Кампањата се состои од \textbf{20 мерни извршувања} --- 10 спарени почетни вредности (42–51) по
два крака. Двата крака за иста почетна вредност работат врз битово идентични податоци, со
идентично поставено труење, на истиот физички јазол, што секоја почетна вредност ја прави
една \textbf{спарена опсервација}.

Редоследот е балансиран: пет почетни вредности се извршени прво со наивниот пристап
(42, 44, 46, 47, 49) и пет прво со SISA (43, 45, 48, 50, 51), со што ефектот на редоследот
не се меша со ефектот на методот. Пред секоја мерна секвенца се чека уредот да достигне
стабилно термичко плато, а страничната кеш-меморија се празни.

Секое извршување ги поминува истите четири фази: обука со активно труење (10 рунди),
означување на компромитираните примери преку оракул-маска, \textbf{обновување} (мерниот прозорец,
изведен локално на јазолот) и повторно приклучување кон федерацијата (5 рунди). Фазата 4 во
двата крака продолжува од \textbf{истиот, номинално отруен} глобален модел, што значи дека
разликите во неа се припишуваат на начинот на обновување, а не на различна почетна точка.

\section{Користени метрики}

Зависните променливи се поделени во две групи, кои водат кон спротивни заклучоци и затоа се
изложени одделно.

\textbf{Физички трошок} (H1–H4, H6), мерен исклучиво во прозорецот на Фаза 3:

\begin{itemize}
  \item \emph{Време до обновување} --- реално време, запишано од самиот модул за обновување.
  \item \emph{Нето-енергија} --- трапезоиден интеграл $\int W\,dt$ врз примероци од мерачот FNB58, со \textbf{одземена основна потрошувачка во мирување}, што ја дава маргиналната цена на обновувањето, а не потрошувачката на целата плоча.
  \item \emph{Термичко оптоварување} --- број секунди со активен бит за придушување. Се читаат само \textbf{живите} битови; битовите што бележат дека настанот се случил некогаш остануваат поставени до рестарт и би го преувеличиле придушувањето.
  \item \emph{Влезно-излезни операции} --- запишани мегабајти на SD-картичката и средно време на чекање за влез/излез.
  \item \emph{Режиски трошок} --- контролни точки запишани во текот на Фаза 1 (само за SISA).
\end{itemize}

\textbf{Корисност} (H5) се оценува врз издвоеното глобално тест-множество од 100\,000 примери, од
кои околу 3\,494 се бенигни. Како што беше образложено во \hyperref[chap:theory]{Дел~\ref*{chap:theory}}, §3.5, одѕивот и
$F_1$ врз такво множество се тривијални, па заклучоците се засноваат на \textbf{специфичноста,
избалансираната точност и MCC}, со ROC-AUC како показател независен од прагот.

\begin{quote}
\textbf{Разграничување што мора да се направи.} Редовите \texttt{p4\_final\_f1} и \texttt{p4\_final\_recall} во
\texttt{paired\_stats.csv} потекнуваат од оценувањето што клиентите го вршат врз сопствената
локална валидациска поделба во текот на рундите. Тие се задржани заради потполност, но
\textbf{не се основа за заклучок}. Секое тврдење за корисноста во §7.5 се потпира на
повторната евалуација на зачуваните глобални модели врз глобалното тест-множество
(\texttt{analysis/evaluate\_utility.py} → \texttt{utility\_evaluation.csv}), со целосна матрица на конфузија.
\end{quote}

Статистиката е спарена по почетна вредност и непараметриска: Вилкоксонов тест на рангирани
знаци, Cliff's $\delta$ како големина на ефектот и bootstrap 95\,\% интервал на доверба на
медијаната на спарената разлика. При $N = 10$ двостраната \textbf{долна граница} на $p$ изнесува
$0{,}0020$; ниту еден тест не може да даде помала вредност, па секој показател што ја
достигнува таа вредност е на максималната достижна значајност. Поради малиот примерок,
тврдењето го носат големината на ефектот и интервалот на доверба, а не самата $p$-вредност.

\section{Анализа на примарни резултати (H1–H3)}

Табелата \ref{tab:paired-cost} ги дава спарените статистики за физичките трошоци.

\input{tables/tab_paired_cost.tex}

Резултатот е недвосмислен и конзистентен низ трите примарни показатели:

\begin{table}[htbp]
\centering
\caption{Спарени статистики за примарните физички показатели (H1--H3), $N = 10$.}
\label{tab:primary-results}
\begin{tabular}{lccccc}
\toprule
Показател & Наивно (медијана) & SISA (медијана) & Однос & $p$ & $\delta$ \\
\midrule
Време до обновување (H1) & 242,1 s & 29,5 s & \textbf{8,2$\times$} & 0,0020 & $+1{,}00$ \\
Нето-енергија (H2) & 0,093 Wh & 0,017 Wh & \textbf{5,4$\times$} & 0,0020 & $+1{,}00$ \\
Секунди со придушување (H3) & 233 s & 27,5 s & \textbf{8,5$\times$} & 0,0020 & $+1{,}00$ \\
\bottomrule
\end{tabular}
\end{table}

\textbf{Cliff's $\delta = +1{,}00$ на сите три показатели значи целосно раздвојување:} ниту едно
извршување со SISA не се преклопува со ниту едно наивно извршување. Опсезите го покажуваат
тоа директно --- времето до обновување е 230,0–294,4 s за наивниот пристап наспроти
26,3–43,8 s за SISA; нето-енергијата 0,079–0,121 Wh наспроти 0,012–0,019 Wh; придушувањето
204–293 s наспроти 2–49 s. Bootstrap интервалот на доверба за разликата во времето
[208,7; 219,6] s не ја содржи нулата.

\textbf{Забелешка за заокружувањето.} Односот на енергијата пресметан од незаокружените
медијани (0,0931 наспроти 0,0173 Wh) изнесува $5{,}4\times$; истата пресметка врз
заокружените вредности од табелата дава $5{,}5\times$. Во текстот се користи вредноста од
незаокружените податоци.

\subsection{Колку од забрзувањето е избегната пресметка?}

Најостриот приговор на овој резултат е дека SISA е побрза затоа што \emph{прави помалку работа},
па споредбата е кружна. Приговорот заслужува броен, а не реторички одговор, па извршената
работа е пресметана директно од манифестите, изразена во \textbf{примерок-ажурирања} (број редови
× број поминувања низ нив):

\begin{itemize}
  \item наивниот крак поминува 10 епохи врз задржаните ≈ 736\,900 реда, односно ≈ 7,37 милиони;
  \item SISA повторно обработува 47 сегменти, од кои --- бидејќи отруените сегменти 3 и 4 од шардот 1 се повторуваат во \textbf{секоја} рунда --- 27 се полни (по 32\,000 реда), а 20 се речиси испразнети (по чистењето во нив остануваат вкупно околу 940 реда). Тоа дава ≈ 873\,000.
\end{itemize}

\begin{table}[htbp]
\centering
\caption{Нормализација на забрзувањето по извршена работа.}
\label{tab:work-normalized}
\begin{tabular}{lc}
\toprule
Величина (медијана, $N = 10$) & Вредност \\
\midrule
Однос на извршена работа (наивно : SISA) & \textbf{8,44$\times$} \\
Однос на измерено време & \textbf{8,23$\times$} \\
Однос на пропусност (SISA : наивно) & \textbf{0,98$\times$} \\
\bottomrule
\end{tabular}
\end{table}

Заклучокот е јасен и треба да се изнесе директно: \textbf{забрзувањето е речиси точно еднакво на
избегнатата пресметка.} Двата крака обработуваат редови со практично иста брзина (0,98×);
SISA не е побрза по ред, туку обработува осумкратно помалку редови за да ја постигне истата
гаранција.

Тоа \textbf{не} ја обезвреднува предноста --- напротив, таа станува \emph{предвидлива}. Забрзувањето не
е случајна последица на имплементацијата, туку структурно својство: тоа се определува од
бројот шардови $S$ и од положбата на компромитираниот сегмент, па може да се пресмета
однапред за дадена конфигурација. Тоа е и единствената чесна формулација на одговорот на
приговорот: \emph{да, SISA прави помалку работа --- тоа е самиот механизам; смислата на споредбата
е што двата крака при тоа постигнуваат иста гаранција за заборавање.}

Мал остаток од разликата сепак не се објаснува со работата, и тој има две спротивставени
компоненти. Наивниот крак плаќа \textbf{термички данок}: неговите епохи се продолжуваат од
17–20 s на 24–25 s како што уредот се вжештува (медијана +20\,\%, до +39\,\%), што низ десетте
епохи троши медијана \textbf{15\,\%} од неговото време. SISA, пак, плаќа \textbf{режија по сегмент} --- нов оптимизатор, нов вчитувач и запишување контролна точка на секои 32\,000 реда наместо на
секои 737\,000. Двете приближно се поништуваат, што е причината пропусноста да излезе 0,98×.

\subsection{Зошто траењето, а не температурата}

Термичкиот резултат бара внимателна формулација, бидејќи наивната интуиција --- „потешкото
оптоварување загрева повеќе“ --- овде не е показателот што раздвојува.

Врвната температура во прозорецот на обновувањето е повисока за наивниот пристап
(медијана 88,7 °C наспроти 86,5 °C), но \textbf{опсезите се преклопуваат}: 87,3–91,7 °C наспроти
82,9–88,9 °C, при $\delta = +0{,}79$. Ефектот е голем, но не е целосно раздвојување. Причината
е физичка: уред без активно ладење под секое одржливо оптоварување се заситува кон истата
горна граница; температурата кажува \emph{дека} уредот е топол, не \emph{колку долго} бил
ограничуван.

Показателот што навистина раздвојува е \textbf{траењето на придушувањето}, со целосно
раздвојување ($\delta = +1{,}00$). Механизмот е јасен од дневникот по епохи на наивниот крак:
првата епоха трае околу 17–20 s, а од петтата натаму се стабилизира на 24–25 s. Тоа не е
варијација на мерењето, туку самото придушување --- уредот се вжештува, фреквенцијата се
ограничува, и секоја следна епоха трае подолго. SISA завршува за околу половина минута,
пред уредот воопшто да се загрее до тој режим.

\subsection{Нулти наод: минималниот такт не раздвојува}

Хипотезата H3 првично предвидуваше и \textbf{деградација на тактот} како показател. Кампањата не
ја поткрепува: медијаната на минималната фреквенција е 2400 MHz во \textbf{двата} крака,
$p = 0{,}25$, $\delta = -0{,}22$, а интервалот на доверба $[-300; 0]$ MHz ја содржи нулата.
Наивниот крак навистина паѓа на 1500 MHz кај некои почетни вредности (42, 45) и на 1800 MHz
кај други (50), но не кај секоја --- паѓањето бара уредот да се вжешти доволно во рамките на
прозорецот, што зависи и од собната температура. Нултиот наод се известува како таков, а
минималниот такт се задржува само како описен показател по извршување.

\section{Анализа на варијабилност и параметри (H4, H6)}

\subsection{H4 --- влезно-излезните операции: хипотезата е потврдена, но не во насоката што се очекуваше}

Мерењето покажува дека SISA навистина го \textbf{префрла} трошокот кон складирањето, и тоа со
целосно раздвојување: 12,05 MB наспроти 2,75 MB запишани на SD-картичката во прозорецот на
обновувањето (опсези 11,4–12,7 наспроти 2,7–3,8 MB; $\delta = -1{,}00$, $p = 0{,}0020$).
Односот е околу $4{,}4\times$ во полза на наивниот пристап.

Ова е единствениот показател на кој наивниот пристап победува, и заслужува да се прочита
точно. Насоката на ефектот ја потврдува хипотезата: SISA плаќа со записи. \textbf{Апсолутната
големина, меѓутоа, ја побива нејзината важност.} Дванаесет мегабајти запишани во текот на
половина минута се далеку под прагот на било какво тесно грло --- SD-картичка од класата A1
со редот од 500 случајни запишувања во секунда би се оптоварила само ако записите беа за
редови на големина поголеми. Средното време на чекање за влез/излез го потврдува тоа:
0,39–0,75\,\% за SISA, наспроти 0,01–0,75\,\% за наивниот пристап --- во двата случаја занемарливо.

Причината е големината на моделот. Со 12\,865 параметри, една контролна точка зазема
174\,688 бајти. Хипотезата H4 беше формулирана како \textbf{тестирана, а не претпоставена}, токму
за случајот кога резултатот е занемарлив да може да се пријави искрено, наместо да се
премолчи. Тоа е случајот овде, и тоа поставува јасно ограничување на заклучокот: тврдењето
„SISA го поместува трошокот кон влезно-излезни операции“ е точно по насока, но при оваа
големина на модел тоа поместување \textbf{нема практична последица}. Каде тоа престанува да важи --- при која големина на моделот записите стануваат доминантни --- останува отворено прашање
(\hyperref[chap:future]{Дел~\ref*{chap:future}}).

\subsection{H6 --- вкупен трошок на сопственост}

SISA плаќа непрекинат режиски трошок во текот на целата обука, без оглед на тоа дали
барањето за заборавање воопшто ќе пристигне. Измерено низ сите десет почетни вредности, тој
трошок е крајно доследен: \textbf{250 контролни точки, 43\,672\,000 бајти (43,67 MB) и 3,1–5,9 s
време за запишување} по извршување, со медијана од 5,56 s.

Спротивставено на тоа, медијаната на заштеденото време при обновувањето изнесува
\textbf{213,7 s} (опсег низ почетните вредности 200,8–262,3 s).

Билансот е јасен: режискиот трошок од околу 5,6 s, распределен низ десет рунди на обука, се
отплаќа \textbf{над триесет пати} со едно единствено обновување. Дури и ако барањето за
заборавање никогаш не пристигне, апсолутната цена останува околу шест секунди и 44 MB по
целосна кампања на обука --- занемарливо. При оваа скала на моделот, SISA е \textbf{евтина
полиса за осигурување}. Ова е графички прикажано на Слика \ref{fig:h6}.

Треба да се наведе и што овој број \textbf{не} покажува: 43,67 MB се однесуваат на десет рунди
обука со мал модел. Трошокот расте линеарно со бројот на рунди, со $S \times R$ и со
големината на моделот, па за поголеми модели или подолги кампањи билансот мора да се
пресмета одново.

\subsection{Варијабилност меѓу почетните вредности}

Времето до обновување кај SISA е изразито репродуцибилно (26,3–43,8 s), но не е константно.
Двете највисоки вредности --- почетните вредности 50 и 51, по околу 43 s --- имаат конкретна и
проверлива причина, која \textbf{не е} термичка: кај тие две извршувања компромитираното множество
е приближно половина од она кај останатите осум (31\,517 и 31\,561 наспроти 63\,060–63\,155
примери). Помалку отстранети редови значи дека во сегментите 3 и 4 остануваат повеќе редови
за повторна обработка, па извршената работа расте од ≈ 873\,000 на ≈ 1\,189\,000
примерок-ажурирања, односно за 36\,\%. Односот работа : време останува стабилен и за нив
(6,46× наспроти 6,10×), што потврдува дека станува збор за повеќе работа, а не за побавна
машина. Причината за различната големина на множеството е разгледана во
\hyperref[chap:limitations]{Дел~\ref*{chap:limitations}}, §8.12.

\subsection{Предноста на SISA зависи од локалноста на компромитирањето}

Бидејќи забрзувањето е точно толку колку што е избегнатата пресметка, следува дека тоа
\textbf{исчезнува кога компромитирањето престане да биде локално}. Тоа не е теоретска
претпоставка --- измерено е.

При едно извршување со несовпаднат кеш, фиксираните компромитирани редови паднаа во
\textbf{сегментот 0 на два шарда} под распоредувањето на тековната почетна вредност, наместо во
сегментите 3–4 на еден шард. Бидејќи не постоеше ниту една чиста контролна точка пред нив,
постапката мораше да ги обработи повторно сите 250 сегменти, што траеше \textbf{264 s} --- практично еднакво на наивниот крак (медијана 242,1 s). Предноста од 8× падна на околу 0,9×.

Ова дава јасна \textbf{граница на важење} на резултатот, изразена во истите единици како и самиот
резултат: очекуваното забрзување е приближно

\[
\frac{S \cdot R \cdot T}{\,|\{\text{сегменти по компромитираниот}\}|\,},
\]

каде бројителот е целосната повторна обука. Моделот на закана „извор компромитиран во
моментот $\tau$“ --- доцни сегменти на еден шард --- е токму оној за кој SISA е конструирана и
во кој нејзината предност е најголема. Компромитирање во ран сегмент, или распространето низ
сите шардови, ја сведува на нула. Ова се известува како \textbf{услов на важење и наод за
скалирањето}, не како недостаток: истата пресметка му овозможува на проектантот однапред да
процени дали SISA воопшто се исплати за неговиот модел на закана.

(Извршувањето со несовпаднат кеш е една неповторена опсервација, настаната ненамерно, и
затоа не влегува во спарената статистика; се известува како илустрација на границата, со
измерената вредност наведена.)

\section{Компаративна анализа --- корисност (H5)}

Овој дел ја носи најважната квалификација на целиот труд. Ако се чита само табелата со
физички трошоци, заклучокот би бил дека SISA доминира. Повторната евалуација на корисноста
покажува дека тоа не е точно.

Табелата \ref{tab:utility} ги дава медијаните по пристап врз глобалното тест-множество по
обновувањето и повторното приклучување.

\input{tables/tab_utility.tex}

\begin{table}[htbp]
\centering
\caption{Корисност при подразбираниот праг $0{,}5$ (Фаза 4, медијани).}
\label{tab:utility-default-threshold}
\begin{tabular}{lcccc}
\toprule
Показател (Фаза 4) & Наивно & SISA & Мнозинска класа & $\delta$ \\
\midrule
Одѕив & 0,960 & 1,000 & 1,000 & $-1{,}00$ \\
\textbf{Специфичност} & \textbf{0,237} & \textbf{0,000} & 0,000 & $+1{,}00$ \\
\textbf{Избалансирана точност} & \textbf{0,598} & \textbf{0,500} & 0,500 & $+1{,}00$ \\
\textbf{MCC} & \textbf{0,177} & \textbf{0,017} & 0,000 & $+1{,}00$ \\
Удел предвидени позитивни & 0,953 & 1,000 & 1,000 & $-1{,}00$ \\
ROC-AUC & 0,725 & 0,751 & 0,500 & $-0{,}42$ \\
\bottomrule
\end{tabular}
\end{table}

Прочитано преку одѕивот, SISA изгледа совршено: 1,000 наспроти 0,960. Прочитано преку
матрицата на конфузија, тој одѕив е \textbf{тривијален}. Обновениот глобален модел со SISA
предвидува дека \textbf{секој} мрежен тек е напад: уделот на предвидени позитивни е точно
1,0000, а бројот на точно детектирани бенигни примери се движи меѓу \textbf{0 и 2 од околу
3\,494}. Избалансираната точност изнесува 0,500 --- вредност на случаен избор --- а MCC е
0,017, практично нула.

Наивниот пристап не е добар во апсолутна смисла, но \textbf{навистина дискриминира}: тој
препознава меѓу 394 и 1\,206 бенигни текови, со избалансирана точност 0,598 и MCC 0,177.
Раздвојувањето е целосно ($\delta = +1{,}00$) на сите три показатели отпорни на нерамнотежа.

Заклучокот, формулиран против сопствениот предмет на истражувањето: \textbf{при прагот на одлука
што реално се распоредува, наивното повторно тренирање е подобро, а „совршениот одѕив“ на
SISA е гласање на мнозинството.}

\subsection{Дали тоа е загубена информација или погрешно поставен праг?}

Прагот од 0,5 не е својство на моделот, туку конвенција. Бидејќи ROC-AUC на SISA (0,751) е
нешто повисок од оној на наивниот пристап (0,725), се поставува прашањето колку од колапсот
е само неправилно поставена граница на одлучување. Тоа се проверува со премин низ сите
прагови: за секој зачуван глобален модел се пресметува највисоката достижна избалансирана
точност.

\begin{table}[htbp]
\centering
\caption{Корисност при подесен (оптимален) праг наспроти подразбираниот праг $0{,}5$.}
\label{tab:utility-tuned-threshold}
\begin{tabular}{lccc}
\toprule
Показател (Фаза 4, медијана) & Наивно & SISA & $\delta$ \\
\midrule
Избалансирана точност при праг 0,5 & 0,598 & 0,500 & $+1{,}00$ \\
\textbf{Избалансирана точност при оптимален праг} & \textbf{0,671} & \textbf{0,709} & \textbf{$-0{,}68$} \\
Специфичност при оптимален праг & 0,662 & 0,694 & $-0{,}20$ \\
MCC при оптимален праг & 0,142 & 0,171 & $-0{,}62$ \\
Оптимален праг & 0,986 & 0,974 & $+0{,}62$ \\
\bottomrule
\end{tabular}
\end{table}

\textbf{Резултатот го превртува заклучокот од претходниот пасус.} Калибрацијата му носи на SISA
+0,209 избалансирана точност, наспроти +0,073 за наивниот пристап; по неа \textbf{SISA е подобра во
9 од 10 почетни вредности} (0,709 наспроти 0,671), и тоа истовремено по специфичност и по
MCC. Колапсот при прагот 0,5 значи \textbf{не е изгубена информација --- тој е чиста калибрација}.

Два детаља што не смеат да се премолчат. Прво, \textbf{двата} модела бараат екстремен праг
(0,96–0,99), што значи дека и двата се лошо калибрирани; SISA е само полошо од наивниот.
Второ --- и поважно --- прагот овде е избран врз \textbf{истото} множество врз кое се и оценува.
Тоа е \emph{оракул} праг: тој ја дава \textbf{горната граница} на она што калибрацијата може да го
постигне, не проценка за распоредување. Реална проценка бара прагот да се избере на
издвоено валидациско множество и да се оцени на дисјунктно тест-множество
(\hyperref[chap:limitations]{Дел~\ref*{chap:limitations}}, §8.13).

\subsection{Дали дефектот доаѓа од заборавањето?}

Тоа е прашањето што ја определува важноста на целата споредба, и одговорот е \textbf{не}.

Повторната евалуација применета врз моделите од \textbf{Фаза 1} дава избалансирана точност 0,500
и за SISA (опсег 0,5000–0,5007) --- пред воопшто да се изврши какво било заборавање. Колапсот
затоа му претходи на обновувањето. Контролна проверка врз зачуваните составни модели го
локализира и точното место: оценети поединечно, петте составни модели препознаваат меѓу 543
и 996 бенигни текови со избалансирана точност 0,545–0,585, додека параметарската средна
вредност на нивните тежини --- токму она што клиентот го испраќа кон FedAvg --- не препознава
ниту еден. Дефектот настанува при \textbf{самото просечување}, не во изолираните шардови и не во
постапката за заборавање. (Детален излез: \texttt{results/msc/diagnostics/}.)

Ова е одлучувачко за толкувањето на трудот: измерената \textbf{цена на заборавањето} е измерена
коректно, а утврдениот недостаток во корисноста е својство на конкретната \textbf{интеграција} на
SISA во федерираната агрегација --- документираното отстапување во кое клиентот испраќа средна
вредност на тежините наместо ансамбл од предвидувања. Тоа е компонента што може да се замени
без да се допре механизмот на заборавање.

За полнота: наивниот крак ја \textbf{подобрува} корисноста по обновувањето --- избалансираната
точност расте од 0,564 во Фаза 1 на 0,598 во Фаза 4 --- што е конзистентно со тоа дека
отстранувањето на 63\,000 погрешно означени примери му помага на моделот.

\section{Импликации за практична примена}

Резултатите водат кон препорака што е поусловена отколку што почетната хипотеза
претпоставуваше.

\begin{enumerate}
  \item \textbf{Кога заборавањето мора да биде брзо и евтино, SISA испорачува убедливо.} Осумкратно пократко обновување, петкратно помала потрошена енергија и осумкратно помалку време во термичко придушување, со целосно раздвојување врз десет спарени извршувања, не е маргинална предност. За уред што мора да остане оперативен додека се обновува --- или што работи на батерија --- разликата меѓу четири минути и половина минута е разликата меѓу прекин и незабележлива операција.
  \item \textbf{Режискиот трошок не е пречка при оваа скала.} Околу шест секунди и 44 MB по кампања на обука се отплаќаат триесеткратно со едно обновување. Стравувањето дека постојаното запишување контролни точки ќе ја оптовари SD-картичката не се потврди за модел со 13\,000 параметри.
  \item \textbf{SISA не смее да се распореди со подразбираниот праг од 0,5.} Тоа е најнепосредната практична поука. Клиент што го земе моделот како што е добива детектор што сè прогласува за напад --- бескорисен во производство колку и детектор што ништо не открива. Дефектот е тивок: одѕивот и F1 изгледаат одлично додека тој трае.
  \item \textbf{Но поправката е бесплатна, и по неа SISA не губи ништо на корисност.} Со подесен праг SISA достигнува избалансирана точност 0,709 наспроти 0,671 за наивниот пристап, подобра во 9 од 10 почетни вредности, и подобра истовремено по специфичност и по MCC. Калибрацијата не бара повторна обука, ниту го менува ниту еден од измерените физички трошоци. Практично: праг подесен на издвоено множество мора да биде дел од распоредувањето, а трајното решение е агрегација што ги чува составните модели одвоени (ансамбл или дестилација), бидејќи поединечно тие дискриминираат подобро од просекот на нивните тежини.
\end{enumerate}

Со други зборови, измерениот наод не е ниту „SISA е подобра во сè“, ниту „SISA купува брзина
со корисност“, туку: \textbf{SISA го прави заборавањето осумкратно поевтино без загуба во
корисноста --- под услов прагот на одлучување да биде калибриран.} Некалибрирана, таа
испорачува модел што е формално одличен и практично бескорисен, и тоа е предупредувањето што
овој труд го носи во проектантската пракса.

